\begin{claim}{C.15}{Proved under stated hypotheses}{An exact resolving functor removes the colimit-exactness requirement}\label{C15}
Let $\mathcal M$ satisfy (a)--(d) and have all small coproducts. Suppose there is an exact cocontinuous $k$-linear functor
\[
 F:\mathcal M\longrightarrow\Ch_{\geq0}(k)
\]
with a specified identification $H_0F=\varphi$, and with $H_iF(P)=0$ for every projective $P\in\mathcal M$ and $i>0$. Then the canonical functor $\mathcal M[W^{-1}]\to D_{\geq0}(k)$ is an equivalence. We require neither AB4 nor AB5, a small projective generating category, nor a right adjoint to $F$. The existence of this exact cocontinuous resolving functor is the added hypothesis; it has been proved in the corollary setting by Claim~\ref{C6}, and is not inferred from (a)--(d) alone.
\end{claim}
\begin{proof}
Claim~\ref{C2} identifies $F$ with $L\varphi$ on objects of $\mathcal M$. Define
\[
 \mathsf{Cof}=F^{-1}\{\text{degreewise monomorphisms}\},
 \qquad W=F^{-1}\{\text{quasi-isomorphisms}\}.
\]
These cofibrations need not be monomorphisms in $\mathcal M$. That distinction allows us to verify the ABC axioms without assuming exactness of colimits in $\mathcal M$.

All objects are cofibrant, identities and composites are cofibrations, and two-out-of-three holds for $W$. The category has all small colimits, since an abelian category with small coproducts has cokernels and hence all small colimits. Pushouts of cofibrations, and pushouts of acyclic cofibrations, have the required properties because $F$ preserves pushouts and the corresponding assertions hold for monomorphisms of complexes of $k$-modules. For a map $f:X\to Y$, the factorization
\[
 X\xrightarrow{(f,\alpha\otimes X)}Y\oplus(J\otimes X)
 \xrightarrow{\mathrm{pr}_Y}Y
\]
has first map in $\mathsf{Cof}$: it is monic in $\mathcal M$ and $F$ is exact. The second map is in $W$ by (d). This proves the factorization axiom. Arbitrary coproducts of cofibrations and of acyclic cofibrations have the required properties after applying $F$, because $F$ preserves coproducts and sums of $k$-modules are exact. Finally, apply $F$ to a countable sequence of cofibrations. Its image is a sequence of degreewise monomorphisms in $\Ch_{\geq0}(k)$; its colimit is again monic from every stage. If the successive cofibrations are acyclic, those maps to the colimit are quasi-isomorphisms, since filtered colimits of $k$-modules are exact. Preservation of the colimit by $F$ proves both parts of the final axiom in $\mathcal M$. Thus $\mathcal M$ is an ABC cofibration category, and $F$ preserves all colimits and cofibrations used to construct its homotopy colimits.

The direct-replacement/Reedy construction in Claim~\ref{C13} therefore proves that $F$ preserves homotopy colimits over every small indexing category. Explicitly, a restricted Reedy cofibrant replacement over $\Delta'I$ remains restricted Reedy cofibrant under $F$, and $F$ takes its colimit to the colimit of its image. Brown's calculus of fractions also shows that the induced $\overline F:\mathcal M[W^{-1}]\to D_{\geq0}(k)$ is conservative: representing a morphism by $X\to Z\leftarrow Y$, with second map in $W$, shows that invertibility of its image forces the first map into $W$ as well. These are the applications of \citet[Theorems 6.4.4 and 9.6.3]{RadulescuBanu2009}; no combinatorial-model recognition theorem is used.

Let $H=\operatorname{hocolim}_{\Delta^{\mathrm{op}}}$ on the localization of simplicial objects at objectwise $W$. Its adjunction to the constant diagram functor has invertible counit, since the nerve of $\Delta^{\mathrm{op}}$ is contractible. It therefore identifies $\mathcal M[W^{-1}]$ with the further localization of simplicial objects at all maps inverted by $H$. For $K\in\Ch_{\geq0}(\mathcal M)$, exactness and additivity of $F$ give
\[
 \overline F(H\Gamma K)
 \simeq\operatorname{hocolim}_{\Delta^{\mathrm{op}}}F(\Gamma K)
 \simeq\Tot\bigl(F_q(K_p)\bigr).
\]
The second comparison is the normalized alternating-face computation of the homotopy colimit in nonnegative complexes; all total degrees have finite diagonals. Conservativity now identifies the maps inverted by $H$ with the maps inverted by $\Tot F$. Exactness of each $F_q$ ensures that ordinary quasi-isomorphisms in $\Ch_{\geq0}(\mathcal M)$ belong to this class. Resolving a complex by a first-quadrant double complex of projectives and using Claim~\ref{C2} identifies $\Tot F$ with $L\varphi$ on $D_{\geq0}(\mathcal M)$. Hence
\[
 \mathcal M[W^{-1}]
 \simeq D_{\geq0}(\mathcal M)
       [(L\varphi)^{-1}(\mathrm{Iso})^{-1}].
\]

It remains to check the canonical counit of $L\varphi\dashv\kappa$. Normalization (c) proves it on $k[0]$. Both $F$ and $\kappa$ preserve coproducts, so it holds on every free $k$-module in degree zero, and hence on every projective one by retracts. Notice that we use preservation by the exact functor $F$, rather than assuming that sums of projective resolutions are exact in $\mathcal M$. Resolve an arbitrary nonnegative $k$-complex by a complex of projective $k$-modules. Exactness of $\kappa$ and the same first-quadrant comparison show that its canonical derived counit is invertible. The localized adjunction therefore gives the equivalence with $D_{\geq0}(k)$, proving the claim.
\end{proof}

